\documentclass[12pt,a4paper]{article}
\usepackage[utf8]{inputenc}
\usepackage[T1]{fontenc}
\usepackage[spanish]{babel}
\usepackage{amsmath,amssymb,amsfonts}
\usepackage{geometry}
\usepackage{booktabs}
\usepackage{hyperref}
\geometry{margin=2.5cm}

\title{Los Horizontes Causales como Andamiaje Conceptual en la F\'isica Gravitacional y Cosmol\'ogica}
\author{Kevin Mu\~noz}
\date{}

\begin{document}
\maketitle

\begin{abstract}
Los desarrollos recientes en física gravitacional, cosmología y teoría cuántica de campos enfatizan cada vez más la significancia de los horizontes causales, la accesibilidad dependiente del observador y la entropía asociada al borde. De manera independiente, los enfoques centrados en el horizonte han explorado la posibilidad de que la estructura causal pueda desempeñar un papel organizador en la interpretación de los fenómenos cosmológicos.

El presente trabajo introduce un andamiaje conceptual mínimo destinado a clarificar las correspondencias estructurales entre varias direcciones de investigación en las que los horizontes ya poseen relevancia física explícita. En lugar de proponer nuevos principios dinámicos, modificar teorías existentes o introducir un marco unificado, el andamiaje identifica un pequeño conjunto de características estructurales recurrentes: primacía causal, dominios definidos por el horizonte y entropía asociada al borde.

Se muestra que estos principios exhiben fuerte compatibilidad conceptual con la termodinámica de horizontes, las perspectivas holográficas generales, la física cosmológica de horizontes y enfoques seleccionados basados en estructura causal.

El andamiaje no debe interpretarse como una teoría o síntesis. Su propósito se limita a identificar una interfaz conceptual común a través de la cual enfoques de otro modo distintos pueden discutirse sin modificación de sus supuestos internos o estructuras formales.
\end{abstract}

\section{Introducción}

Los horizontes causales aparecen en una amplia variedad de contextos gravitacionales y cosmológicos. Los horizontes de agujeros negros, los horizontes cosmológicos y las fronteras causales dependientes del observador delimitan consistentemente la accesibilidad mientras exhiben simultáneamente propiedades termodinámicas e informacionales.

A pesar de su aparición recurrente, los horizontes se tratan con frecuencia dentro de contextos específicos de cada marco, en lugar de como estructuras organizadoras potencialmente compartidas.

El presente trabajo no busca unificar teorías existentes ni introducir nuevos principios físicos. En cambio, explora si un pequeño conjunto de ideas estructurales asociadas a los horizontes puede funcionar como una interfaz conceptual mínima entre enfoques de otro modo distintos.

El objetivo es, por tanto, organizacional en lugar de explicativo.

El marco desarrollado aquí debe entenderse como un andamiaje conceptual destinado a facilitar el diálogo y la comparación mientras preserva la independencia teórica de los enfoques participantes.

\section{Estatus y Alcance de las Afirmaciones}

El andamiaje propuesto en este trabajo es de carácter conceptual.

No introduce:

\begin{itemize}
  \item nuevos principios dinámicos
  \item modificaciones de la Relatividad General
  \item descripciones microscópicas
  \item supuestos físicos adicionales
  \item predicciones observacionales
\end{itemize}

La compatibilidad dentro del andamiaje no debe interpretarse como:

\begin{itemize}
  \item equivalencia
  \item reducción
  \item unificación
  \item ontología compartida
\end{itemize}

En cambio, la compatibilidad se refiere únicamente a la presencia de elementos estructurales recurrentes.

Las afirmaciones relativas a correspondencia deben, por tanto, interpretarse como afirmaciones conceptuales en lugar de físicas.

No se realiza ningún intento de derivar un marco a partir de otro.

\section{Criterios de Compatibilidad Conceptual}

El propósito de esta sección es identificar criterios bajo los cuales la correspondencia estructural puede discutirse significativamente.

Estos criterios no funcionan como requisitos de validación y no deben interpretarse como restricciones a la construcción de teorías.

En cambio, especifican las condiciones mínimas para la inclusión dentro del presente andamiaje.

\subsection{Primacía Causal}

El primer criterio concierne el papel de la estructura causal.

Los enfoques considerados compatibles dentro del andamiaje asignan significancia fundamental a la accesibilidad causal o tratan las relaciones causales como estructuralmente previas a la descripción puramente volumétrica.

No se requieren supuestos respecto al estatus ontológico.

\subsection{Dominios Definidos por el Horizonte}

Un segundo criterio concierne el papel de los horizontes causales.

Los enfoques compatibles reconocen los horizontes como estructuras que delimitan dominios físicamente significativos de accesibilidad y relevancia informacional.

Los horizontes no necesitan poseer realización matemática idéntica en distintos marcos.

El requisito es estructural en lugar de formal.

\subsection{Entropía Asociada al Borde}

Un tercer criterio concierne la entropía.

Los enfoques incluidos dentro del andamiaje asignan significancia física a la entropía asociada a fronteras causales, en lugar de exclusivamente a los grados de libertad volumétricos encerrados.

Este criterio refleja los conocimientos recurrentes de la termodinámica gravitacional y la física de horizontes.

No se requieren supuestos respecto al origen microscópico.

\subsection{Alcance y Exclusiones}

La presencia de los criterios anteriores no implica acuerdo respecto a:

\begin{itemize}
  \item estructura microscópica
  \item ontología del espacio-tiempo
  \item mecanismos de gravedad cuántica
  \item ecuaciones dinámicas
  \item completación global
\end{itemize}

Los enfoques basados fundamentalmente en puntos de vista externos privilegiados o descripciones puramente volumétricas pueden caer fuera del alcance previsto.

No se implica ninguna afirmación de incompletitud o inadecuación.

\section{Termodinámica de Horizontes}

Las propiedades termodinámicas asociadas a los horizontes proporcionan uno de los ejemplos más claros de compatibilidad estructural.

La termodinámica de agujeros negros demuestra que la entropía y la temperatura surgen naturalmente en asociación con fronteras causales.

Importantemente, tales propiedades dependen de la estructura del horizonte en lugar de descripciones detalladas de regiones inaccesibles.

Desde la perspectiva adoptada aquí, la termodinámica de horizontes ilustra cómo las fronteras causales pueden restringir las descripciones físicas a través de consideraciones estructurales.

El andamiaje no intenta derivar la dinámica gravitacional a partir de relaciones termodinámicas.

Sin embargo, enfoques existentes sugieren que las formulaciones efectivas basadas en termodinámica de horizontes pueden reproducir consistentemente el comportamiento gravitacional familiar.

Tales desarrollos motivan, pero no requieren, la presente perspectiva organizacional.

Como ejemplo concreto dentro del presente programa, la primera ley $dE = TdS - PdV$ aplicada al horizonte aparente en $R_A = 1/H$ reproduce la ecuación de Friedmann estándar, demostrando que la dinámica cosmológica puede reformularse como una afirmación termodinámica sobre la frontera causal.

\section{Perspectivas Holográficas Generales}

Las ideas holográficas generales asocian frecuentemente la información físicamente relevante con estructuras de borde.

Cuando se interpretan ampliamente, tales perspectivas exhiben superposición conceptual con descripciones centradas en el horizonte en las que las condiciones de borde restringen descripciones efectivas del bulk.

El presente andamiaje no asume la validez de dualidades holográficas específicas.

Tampoco requiere correspondencias exactas borde-bulk.

La relevancia de las perspectivas holográficas reside únicamente en su énfasis estructural en la información definida por el borde.

La compatibilidad debe, por tanto, entenderse como limitada y organizacional.

\section{Horizontes Cosmológicos en Espacio-Tiempos en Expansión}

Los horizontes cosmológicos delimitan naturalmente la accesibilidad observacional dentro de espacio-tiempos en expansión.

Estas estructuras exhiben propiedades termodinámicas análogas a las encontradas en sistemas de agujeros negros y, por tanto, satisfacen varios de los criterios introducidos anteriormente.

Dentro de los enfoques centrados en el horizonte, los dominios cosmológicos observables pueden interpretarse como regiones causalmente acotadas.

Tales interpretaciones permanecen plenamente compatibles con la fenomenología cosmológica convencional.

No se propone ninguna dinámica cosmológica alternativa.

El andamiaje simplemente señala que los horizontes causales proporcionan un lenguaje estructural común a través del cual tales descripciones pueden discutirse.

En la realización efectiva, las perturbaciones del horizonte $\psi = \delta R_A / R_A$ proporcionan una descripción basada en el borde de las perturbaciones de densidad mediante $\delta\rho/\rho = -2\psi$, ilustrando concretamente cómo las fluctuaciones del borde codifican los observables del bulk.

\section{Estructura Causal y Geometría Emergente}

Varios enfoques asignan significancia fundamental a las relaciones causales mientras tratan la geometría como emergente o efectiva.

Estas perspectivas frecuentemente definen dominios físicamente significativos a través de relaciones de accesibilidad en lugar de a través de la inclusión en estructuras espaciales preexistentes.

Tales características exhiben fuerte correspondencia estructural con las descripciones centradas en el horizonte.

La compatibilidad no implica ontología compartida, realización matemática ni mecanismo físico.

El andamiaje identifica únicamente patrones recurrentes de organización conceptual.

\section{Límites e Invitaciones Abiertas}

El presente andamiaje es intencionalmente mínimo.

No proporciona:

\begin{itemize}
  \item mecanismos explicativos
  \item derivaciones físicas
  \item estructura predictiva
  \item síntesis formal
\end{itemize}

Su papel se restringe a identificar superposición conceptual.

Varias preguntas permanecen fuera del alcance, incluyendo:

\begin{itemize}
  \item origen microscópico de los grados de libertad del horizonte
  \item mapeos borde-a-bulk
  \item realizaciones dinámicas efectivas
  \item completación gravitacional cuántica
\end{itemize}

Los desarrollos futuros pueden aclarar si existen formas más fuertes de correspondencia.

Tales preguntas permanecen abiertas.

\section{Conclusión}

Hemos introducido un andamiaje conceptual centrado en el horizonte destinado a identificar principios estructurales recurrentes entre enfoques gravitacionales y cosmológicos en los que los horizontes causales desempeñan papeles significativos.

El andamiaje no debe interpretarse como una teoría o marco unificador.

Su contribución se limita a la identificación de una interfaz conceptual mínima organizada en torno a:

\begin{itemize}
  \item primacía causal
  \item dominios definidos por el horizonte
  \item entropía asociada al borde
\end{itemize}

Al clarificar estas correspondencias estructurales, el andamiaje puede facilitar el diálogo entre programas de investigación de otro modo distintos, preservando su independencia formal.

Si estas correspondencias admiten finalmente una realización física más profunda sigue siendo una pregunta abierta.

\end{document}
